\chapter{Recurrent models and gated state}

Read the bases A, C, G, T one at a time. After the last base arrives,
what information about the first base is still available? A fixed-width
window answers by retaining a specified recent interval. A recurrent model
instead maintains a state that it updates at every step. The state is a
learned representation of history, not an automatically lossless copy of it.

This chapter builds recurrent arithmetic from scratch. We first expose the
state transition, then add gates, then test what changes at padding and
chunk boundaries. The main implementation is a manually written GRU cell,
checked against PyTorch in both forward values and gradients. A smaller
scalar recurrence gives exact numbers that can be verified without a library.

We assume Chapter 5's embedding shapes and Chapter 4's loss and gradient
conventions. All numerical probes use assigned weights or synthetic inputs.
They establish properties of the computation, not trained genomic capability.

\section{Unroll one state transition}

Let $x_t\in\mathbb R^D$ be the embedding of the base arriving at step $t$,
and let $h_{t-1}\in\mathbb R^H$ be the previous hidden state.
A simple \Term{recurrent neural network} updates
\begin{equation}
h_t=\tanh(W_xx_t+W_hh_{t-1}+b),
\label{eq:rnn}
\end{equation}
where $W_x\in\mathbb R^{H\times D}$, $W_h\in\mathbb R^{H\times H}$,
and $b\in\mathbb R^H$. The hyperbolic tangent acts elementwise.
The same parameters are reused at every step.

\Plate{01-unroll}{A causal recurrence unrolled over three observed bases. Each state uses the current embedding and the previous state; the shared transition parameters do not become separate trainable copies. A readout after observing A predicts the next base, C. Arrows show computation dependencies, not biological reactions.}{fig:unroll}

For a batch, the code stores $x$ as $[B,T,D]$ and each state as $[B,H]$.
The mathematical column-vector multiplication becomes
\texttt{F.linear(x, weight, bias)} for batch rows: PyTorch applies the
transpose of the stored $[H,D]$ matrix to those rows.
Confusing tensor storage with the mathematical map is a common source of
accidentally transposed weights.

For next-base prediction, a readout produces
\begin{equation}
\ell_t=W_yh_t+b_y\in\mathbb R^4,\qquad
p(s_{t+1}\mid s_{1:t})=\operatorname{softmax}(\ell_t).
\end{equation}
The state at step $t$ has already consumed $s_t$. Scoring it against
$s_t$ would let it predict a base present in its own input. A unidirectional
recurrence is causal only relative to correctly aligned inputs and targets.

We begin each independent record with a declared initial state, here zero.
Two adjacent records in a data loader are not adjacent pieces of one genome
unless the data contract says so. Carrying state across unrelated records
creates an undeclared input channel.

\section{Why the gradient can lose the distant past}

For a single example, define
$a_t=W_xx_t+W_hh_{t-1}+b$. Differentiating Eq.~\ref{eq:rnn} gives
\begin{equation}
J_t=\frac{\partial h_t}{\partial h_{t-1}}
=\operatorname{diag}(1-h_t^2)W_h.
\end{equation}
The influence of an earlier state on a later one is a product:
\begin{equation}
\frac{\partial h_T}{\partial h_k}
=J_TJ_{T-1}\cdots J_{k+1}.
\label{eq:jacobian-product}
\end{equation}
Repeated contractions can suppress a signal; expanding directions can
amplify it. Saturated tanh coordinates contribute small local derivatives.
The foundational LSTM paper motivates gated memory through this
long-range gradient problem \cite{lstm-original}.

A scalar linear recurrence makes the product visible without nonlinear
complications. If $h_t=\alpha h_{t-1}$, then
$h_T=\alpha^T h_0$ and $\partial h_T/\partial h_0=\alpha^T$.
For $\alpha=0.5$, ten steps retain $1/1024$ of a perturbation; for
$\alpha=1.1$, amplification grows with the number of steps.
Increasing a learning rate does not restore information removed by the
forward state or turn a badly conditioned Jacobian product into a stable one.

Shared weights also receive contributions from multiple time steps.
A gradient with respect to $W_h$ is not merely the final state's local
derivative: each earlier occurrence of the same parameter can affect later
states. Our chunk experiment will expose what happens when some of these
paths are deliberately cut.

\section{Separate retention, writing, and reading}

A modern \Term{long short-term memory} cell carries two vectors,
hidden output $h_t$ and cell state $c_t$, both in $\mathbb R^H$.
Let the input gate $i_t$, forget gate $f_t$, and output gate $o_t$ be
sigmoid outputs, and let $g_t$ be a tanh candidate:
\begin{align}
i_t&=\sigma(W_{xi}x_t+W_{hi}h_{t-1}+b_i),\\
f_t&=\sigma(W_{xf}x_t+W_{hf}h_{t-1}+b_f),\\
g_t&=\tanh(W_{xg}x_t+W_{hg}h_{t-1}+b_g),\\
o_t&=\sigma(W_{xo}x_t+W_{ho}h_{t-1}+b_o).
\end{align}
The state update is
\begin{equation}
c_t=f_t\odot c_{t-1}+i_t\odot g_t,\qquad
h_t=o_t\odot\tanh(c_t).
\label{eq:lstm}
\end{equation}
Here $\odot$ is elementwise multiplication. These are the modern
forget-gate equations used by PyTorch's unprojected cell \cite{lstm-api},
not a claim that every feature of this variant appeared in the 1997 paper.

\Plate{02-lstm}{Separate LSTM state paths. The forget gate scales retained cell contents; the input gate scales a candidate write; their sum updates the cell. The output gate controls the exposed tanh-transformed state. Dashed arrows deliver gate values, while solid arrows carry state or candidate values. This diagram omits peepholes and projections.}{fig:lstm}

Take one coordinate with $c_{t-1}=2$, $f_t=0.75$, $i_t=0.2$,
$g_t=-0.5$, and $o_t=0.6$. Retention contributes 1.5, writing contributes
$-0.1$, and $c_t=1.4$. The output is $0.6\tanh(1.4)\approx0.531$.
The cell contents and the exposed output are not the same number.

Holding gates and candidate fixed, the direct derivative along the cell
path is $\partial c_t/\partial c_{t-1}=f_t$ coordinatewise.
If $f_t$ remains near one and writes are small, that path can preserve
information. But the \emph{total} derivative also includes paths through
$h_{t-1}$ and the gates. A direct retention path is not a theorem that
all gradients remain constant or that a trained model will learn the desired memory.

\Listing{lstm_step}

The code packs four gate blocks in the order input, forget, candidate,
output. It accepts two bias vectors because PyTorch stores input-side and
hidden-side biases separately. Their sum plays the role of each $b$ above.
Changing packing order preserves tensor shapes while changing the model.
The parity tests therefore compare actual values, not just dimensions.

\section{Build a GRU cell and name its convention}

A \Term{gated recurrent unit} keeps a single state vector. In the
PyTorch convention used here, reset and update gates are
\begin{align}
r_t&=\sigma(W_{xr}x_t+b_{xr}+W_{hr}h_{t-1}+b_{hr}),\\
z_t&=\sigma(W_{xz}x_t+b_{xz}+W_{hz}h_{t-1}+b_{hz}),
\end{align}
and the candidate and update are
\begin{align}
n_t&=\tanh\left(W_{xn}x_t+b_{xn}
 +r_t\odot(W_{hn}h_{t-1}+b_{hn})\right),\\
h_t&=(1-z_t)\odot n_t+z_t\odot h_{t-1}.
\label{eq:gru}
\end{align}
Thus a large $z_t$ retains the old state in this convention.
Some presentations attach the symbol $z$ to the opposite mixing weight;
read the equation rather than memorizing the gate name \cite{gru-api}.

\Plate{03-gru}{The implemented GRU candidate and state merge. Reset acts after the hidden projection, including its bias. The update gate chooses coordinatewise retention versus candidate contribution. A separate numerical counterexample shows that resetting before a matrix multiplication generally produces different values.}{fig:gru}

For each coordinate, the final update is a convex combination of the
candidate and previous state because a sigmoid lies between zero and one.
That fact bounds the new coordinate between those two values; it does not
bound every derivative of the state-dependent gates.

\Listing{gru_step}

The implementation splits the input and hidden affine projections into
three blocks ordered reset, update, candidate. The two candidate projections
must remain separate until reset has been applied. Adding all affine
outputs and then splitting would lose the information needed to apply
the reset only to the hidden-side candidate contribution.

\subsection{Reset-before and reset-after are not algebraic rearrangements}

PyTorch documents a difference from the original GRU formulation:
its reset is applied to the projected hidden state, rather than before
that projection \cite{gru-variant}. Generally,
\begin{equation}
r\odot(Wh+b)\ne W(r\odot h)+b.
\end{equation}
For $h=(1,2)$, $r=(0.2,0.8)$, zero bias, and
\begin{equation}
W=\begin{bmatrix}1&1\\1&-1\end{bmatrix},
\end{equation}
reset-after gives $(0.6,-0.8)$ whereas reset-before gives $(1.8,-1.4)$.
Both programs can execute and have identical tensor shapes. Only a
convention-aware numerical test detects the mismatch.

This is why a checkpoint needs more than dimensions and vocabulary identity.
It also needs an architecture definition: gate order, bias convention,
activation functions, and where each gate acts. The name GRU is insufficient
to guarantee parameter interchangeability across implementations.

\section{Verify forward values and backward paths}

Our reference artifact assigns deterministic weights.
It compares our GRU step with the library's \texttt{GRUCell}. The test checks the
output and gradients of a squared-output loss with respect to inputs,
previous state, both weight matrices, and both biases.
The tolerance is $10^{-12}$ in CPU float64; no bitwise cross-hardware
guarantee is implied.

There are three complementary checks. A hand case sets all weights and
biases to zero: both gates become 0.5, the candidate becomes zero, and the
new state is half the old state. Library parity tests the full packed
implementation. Numerical finite-difference gradient checking tests the
derivatives independently of a second handwritten backward formula.

Forward parity alone is not enough. A misplaced detach can preserve every
output while removing gradient paths. Conversely, a finite-difference
check on one scalar output is not evidence that a data pipeline masks
padding correctly. Tests should mirror distinct claims.

For $D=4$ and $H=8$, a cell with two bias vectors has
$g(HD+H^2+2H)$ parameters, where $g$ is the number of packed gate blocks.
That gives 112 for a simple tanh cell, 336 for a GRU, and 448 for an LSTM.
LSTM additionally carries both $h$ and $c$ at runtime. Parameter count is
not a benchmark of speed, accuracy, or useful biological memory.

\section{Measure a controlled retention path}

Freeze a GRU update gate to scalar $z$ and set its candidate to zero.
Then $h_t=zh_{t-1}$ and both the state and its direct sensitivity decay
as $z^t$ when $h_0=1$. A half-life in \Term{token steps} satisfies
\begin{equation}
z^\tau=\tfrac12,\qquad
\tau=\frac{\ln(1/2)}{\ln z},\quad 0<z<1.
\end{equation}
For $z=0.5,0.9,0.99$, half-lives are approximately 1, 6.58, and
68.97 steps. At 100 steps, the retained factors are approximately
$7.89\times10^{-31}$, $2.66\times10^{-5}$, and 0.366.

\Plate{04-retention}{Computed retention powers for fixed gates and zero candidate, on a logarithmic vertical scale. This isolates one memory path, not a trained model's total Jacobian. The companion dependency sketch marks additional state-dependent gate paths that the simplified calculation holds fixed.}{fig:retention}

The time unit is part of the interpretation. A hundred base tokens span
100 bases; a hundred disjoint width-six tokens typically span 600 bases.
An overlapping tokenizer advances differently again. A gate half-life
cannot be reported as a genomic distance without the representation contract.

Nor is a slowly decaying state automatically useful. It can preserve an
irrelevant feature, mix several signals destructively, or fail to expose
stored information to the readout. These probes isolate retention mechanics;
learning useful memory requires a task, data split, objective, and evaluation.

\section{Batch sequences without advancing padded time}

Our scan consumes $x$ of shape $[B,T,D]$ and a valid-length vector.
At a valid step, each row updates normally. Beyond its length, the row
retains its previous state. This is a state-update mask, separate from
the target-loss mask introduced earlier.

\Listing{scan_gru}

The output at padded positions repeats the last valid state; it is not
automatically a zero vector. A downstream loss must still ignore invalid
targets. A zero input is not a substitute for holding the state: recurrent
weights and biases can change the state even when the input vector is zero.

The implementation computes candidates before selecting valid rows.
Its padding tests use finite values. A mask does not generally sanitize NaNs
introduced inside unselected computations, especially during backward
propagation; validate inputs rather than relying on masking to repair them.

An empty time dimension returns an empty output sequence and the unchanged
initial state. A zero-length row in a nonempty batch likewise remains at its
initial state. Lengths must be integer counts between zero and the available
number of steps.

\section{Carry, detach, and reset mean different things}

Consider the scalar recurrence
\begin{equation}
h_t=ah_{t-1}+x_t,\quad
a=0.5,\quad h_0=0,\quad x=(1,2,3,4).
\label{eq:scalar}
\end{equation}
Its successive states are $1$, $2.5$, $4.25$, and $6.125$.
Split the computation after step two. Carrying $h_2=2.5$ into the next
chunk preserves the result. Replacing it with zero does not.

\Plate{05-chunks}{Three chunk-boundary policies. Carry keeps both numerical state and its gradient history. Detach carries the same value but cuts backward paths through the boundary. Reset replaces that value with zero and changes the forward computation. Record boundaries and truncated training boundaries therefore require different policies.}{fig:chunks}

\Listing{chunk_experiment}

With full history, the final state expands to
$h_4=a^3x_1+a^2x_2+ax_3+x_4$. Hence
\begin{equation}
\frac{\partial h_4}{\partial x}=(a^3,a^2,a,1)
=(0.125,0.25,0.5,1)
\end{equation}
and
\begin{equation}
\frac{\partial h_4}{\partial a}
=3a^2x_1+2ax_2+x_3=5.75.
\end{equation}
If we detach $h_2$, the second chunk treats 2.5 as a constant:
$h_4=a(a\cdot2.5+3)+4$. Its derivative with respect to the shared $a$
is $2a\cdot2.5+3=5.5$, and the first two input gradients are zero.
Resetting instead gives $h_4=a\cdot3+4=5.5$, with derivative 3.

\begin{center}\small\input{\Generated/results-table.tex}\end{center}

This is \Term{truncated backpropagation} in its most inspectable form:
detachment limits credit assignment, not necessarily the numerical
history carried forward. It is not equivalent to full-sequence training.
Our demonstration keeps parameters fixed while scanning. If an optimizer
updates weights between chunks, a carried state was computed under earlier
weights; that is another distinction from a single full unroll.

For inference with fixed parameters, correctly carried GRU state produces
the same outputs as an uninterrupted scan in our tests. Across independent
records, reset instead. Across padded steps, hold. A single generic
state-clearing rule cannot implement all three boundaries correctly.

\section{What this chapter builds, and what comes next}

Run the reference probes and tests:
\begin{Verbatim}[fontsize=\small]
python3 drafts/ch06/reference.py
python3 -m pytest tests/test_ch06_reference.py
\end{Verbatim}
The artifact covers manual GRU and modern LSTM cells, state masking,
streaming equivalence, finite-difference checks, fixed-gate retention,
and a detachment counterexample. It does not train or benchmark a genomic
language model. DOGMA remains the non-Transformer DNA-native research
direction, Hermon DNA the Transformer-based direction, and Evolutor the
broader framework.

A sequential scan uses a fixed-size state, but nonlinear state dependence
creates a dependency from one step to the next. Chapter 7 asks which state
transitions can be reorganized into efficient scans and how selective
state-space models change the design tradeoffs. The question is no longer
whether there is memory, but how it is represented, updated, and computed.

\Needspace{10\baselineskip}
\section{Exercises with worked solutions}

\subsection*{1. Align the target}
After consuming A, C, G, which base does the final readout predict?
\textbf{Solution:} the next base after G, not G itself. The state has
already received the G embedding.

\subsection*{2. Derive a scalar gradient}
For $h_t=0.8h_{t-1}$, find $\partial h_{10}/\partial h_0$.
\textbf{Solution:} $0.8^{10}\approx0.1074$. This is exact for the fixed
linear recurrence, not for a learned nonlinear gate.

\subsection*{3. Trace an LSTM coordinate}
Use $c=2$, $f=0.75$, $i=0.2$, $g=-0.5$, $o=0.6$.
\textbf{Solution:} the new cell is $1.5-0.1=1.4$ and the hidden output
is about 0.531. The output gate does not erase the underlying cell by
merely hiding part of its readout.

\subsection*{4. Interpret the update gate}
In Eq.~\ref{eq:gru}, what happens as $z$ approaches one?
\textbf{Solution:} the previous state is retained. Approaching zero
instead chooses the candidate. State the convention when comparing equations.

\subsection*{5. Find a shape-preserving bug}
Swap the reset and update blocks in packed weights.
\textbf{Solution:} tensor dimensions still match, but the equations change.
Forward and backward parity with a nontrivial assigned-weight case should fail.

\subsection*{6. Test reset order}
Evaluate the chapter's $W$, $r$, and $h$ counterexample.
\textbf{Solution:} $Wh=(3,-1)$ gives reset-after $(0.6,-0.8)$;
$r\odot h=(0.2,1.6)$ gives reset-before $(1.8,-1.4)$.

\subsection*{7. Count parameters}
Find the GRU parameter count for $D=4,H=8$ with both biases.
\textbf{Solution:} $3(32+64+16)=336$. Excluding one bias vector would
change the count and potentially the checkpoint convention.

\subsection*{8. Convert a memory horizon}
What base distance corresponds to 50 disjoint full-width six-mer steps?
\textbf{Solution:} 300 bases. Do not reuse that conversion for stride-one
overlapping windows, which advance one base per step.

\subsection*{9. Explain padding drift}
Why can feeding zero embeddings through padded steps change final state?
\textbf{Solution:} recurrent projections and biases remain active.
Hold the row's state explicitly and mask invalid targets separately.

\clearpage
\subsection*{10. Separate value and gradient}
Why do full carry and detachment both yield 6.125 but different gradients?
\textbf{Solution:} detachment preserves the boundary's numerical value
while treating it as constant for backward propagation. It removes paths
from the final loss through the first chunk.

\subsection*{11. Prevent cross-record leakage}
A loader places two unrelated sequences consecutively in the same batch row.
\textbf{Solution:} reset state at the record boundary. Detaching alone
retains the earlier record's numerical information.

\subsection*{12. Design a memory evaluation}
How would you test learned long-range genomic dependencies?
\textbf{Worked rubric:} define the biological target, split related records
appropriately, vary dependency distance, compare local-window baselines,
audit strand and position conventions, and report uncertainty across runs.
Separate synthetic retention probes from learned-task accuracy and from
biological generalization.
